\begin{definition}\label{long}\textbf{Longitud de curva.}
    Sean $C$ una curva simple en $\mathbb{R}^n$ y  $\boldsymbol{\sigma}:I\subseteq\R\to\Rn{n}$ una parametrizaci\'on regular y  de clase $\mathcal{C}^1$ a trozos de $C$  (ver\: \ref{param}). Se define la \emph{longitud de} $C$, y se denota por  Long($C$), a 
    \[
        \text{Long}(C)=\sum_{i=0}^{m-1}\int_{I_i}\|\boldsymbol{\sigma}'(t) \|\,dt ,    
    \]
   donde $ I_i= (t_i, t_{i+1})$,  $\boldsymbol{\sigma}\rvert_{I_i}$ es de clase $\mathcal{C}^1$  y  $\{t_i\}_{i=0}^{m}$  es una partici\'on de $I$.
\end{definition}



\begin{example}
  Calcular la longitud de $C$,  siendo $C$  la curva simple  que describe una circunferencia de radio $r$ centrada en el origen. 
  \end{example}  

% Sabemos que podemos describir a  la curva $C$ como   \[  C = \bigl\{(x,y)\in\Rn{2} \;:\; x^2 + y^2 = r^2 \bigr\},\]

\noindent\textbf{Solución:}  Para calcular la longitud de la curva utilizando la definición~\ref{long}, necesitamos realizar tres pasos fundamentales:
\begin{enumerate}
    \item Hallar una parametrización $\boldsymbol{\sigma}$ adecuada para la  curva. 
    \item Calcular el vector tangente $\boldsymbol{\sigma}'$.  
    \item Obtener la norma $\|\boldsymbol{\sigma}'\|$ y resolver la integral correspondiente. 
\end{enumerate}

\noindent Comenzamos parametrizando la curva. Es sencillo ver que 
        \[
            \boldsymbol{\sigma}(t) 
            = 
            \bigl(r\cos t,\;r\sin t\bigr), \quad
            \text{ con } t\in [\,0,\,2\pi\,].
        \] 
       es una parametrización regular y de clase  \(\mathcal{C}^1\)  de $C$. Luego  
        \begin{align*} 
\boldsymbol{\sigma}'(t) &= (-r\sin t,\, r\cos t), \\[0.2cm]
\|\boldsymbol{\sigma}'(t)\| &= \sqrt{(-r\sin t)^2 + (r\cos t)^2} = r\sqrt{\sin^2 t + \cos^2 t} = r. 
\end{align*} 
       Para calcular la longitud de $C$  basta con plantear y resolver la siguiente integral
        \[
            \text{Long}(C)
            =
            \int_{0}^{2\pi} \bigl\|\boldsymbol{\sigma}'(t)\bigr\| \, dt
            =
            \int_{0}^{2\pi} r \; dt
            =
            2\pi r.
        \]
        Por lo tanto, la  longitud de la curva \(2\pi r\), resultado ya conocido, dado que $C$ es una circunferencia de radio \(r\).  \hfill $\blacksquare$ 
  
    
\begin{example}\label{ejer:param_seg}
  Sea la curva simple $S$ dada por el segmento que une los puntos $\mathbf{p}$ y $\mathbf{q}$ en $\mathbb{R}^n$. Calcular la longitud de $S$. 
\end{example}
    
\noindent\textbf{Solución:} Al igual que en el ejercicio anterior, el procedimiento consiste en parametrizar la curva, hallar su vector tangente y luego integrar su norma. 

Comenzamos parametrizando el segmento que va desde el punto $\mathbf{p}$ hacia el punto $\mathbf{q}$.  
\[ 
\boldsymbol{\sigma}(t) = \mathbf{p} + t(\mathbf{q} - \mathbf{p}), \quad  \text{ con } t\in [0,1]. 
\]
Esta parametrización es de clase $\mathcal{C}^1$ y regular de $S$.  Al calcular su vector tangente obtenemos un vector constante:
\begin{align*} 
\boldsymbol{\sigma}'(t) &=\mathbf{q} - \mathbf{p}, \\[0.2cm]
\|\boldsymbol{\sigma}'(t)\| &=   \|\mathbf{q} - \mathbf{p}\|. 
\end{align*} 
Para calcular la longitud de $S$  basta con plantear y resolver la siguiente integral
\[ 
\text{Long}(S) = \int_{0}^{1} \bigl\|\boldsymbol{\sigma}'(t)\bigr\| \, dt = \int_{0}^{1} \|\mathbf{q} - \mathbf{p}\| \, dt = \|\mathbf{q} - \mathbf{p}\|.
\] 
Por lo tanto, la longitud del segmento es $\|\mathbf{q} - \mathbf{p}\|$, resultado que coincide perfectamente con la distancia euclídea elemental entre dos puntos.  \hfill $\blacksquare$ 


\begin{tarea} Calcular, usando la definici\'on  \ref{long},   el per\'imetro de un rombo cuyas diagonales menor y mayor miden $d$ y $D$ respectivamente.  
 \end{tarea}
